Data kan ook aangeleverd worden als output van en ander commando. De data kan dan via het pipe-character (|) doorgestuurd worden naar ons script.

Een Python script dat data van een pipe aanneemt en regel voor regel weergeeft op het scherm zou er zo uit kunnen zien:
\lstinputlisting[language=Python]{code/pipe-json.py}
Zoals je ziet hebben we voor de script de sys-module nodig.

Een commando om dit script te testen op een Linux systeem zou er zo uit kunnen zien:
\begin{lstlisting}[language=bash]
echo "echo '{"Naam":"Jason","Leeftijd":42}' | python pipe-json.py
\end{lstlisting}

We kunnen de input combineren met het werken met opties, zodat je in een optie de key kan meegeven die je wilt zien:
\lstinputlisting[language=Python]{code/pipe-json-arg.py}


